%----------------------------------------------------------------------------------------
\documentclass[a4paper]{article} % Uses article class in A4 format

%----------------------------------------------------------------------------------------
%	FORMATTING
%----------------------------------------------------------------------------------------

\addtolength{\hoffset}{-2.25cm}
\addtolength{\textwidth}{4.5cm}
\addtolength{\voffset}{-3.25cm}
\addtolength{\textheight}{5cm}
\setlength{\parskip}{0pt}
\setlength{\parindent}{0in}

%----------------------------------------------------------------------------------------
%	PACKAGES AND OTHER DOCUMENT CONFIGURATIONS
%----------------------------------------------------------------------------------------

\usepackage[utf8]{inputenc} % Use UTF-8 encoding
\usepackage[T1]{fontenc}
\usepackage{lmodern}
\usepackage[english, spanish]{babel} % Language hyphenation and typographical rules

\usepackage{blindtext} % Package to generate dummy text
\usepackage{charter} % Use the Charter font
\usepackage{microtype} % Slightly tweak font spacing for aesthetics

\usepackage{amsthm, amsmath, amssymb} % Mathematical typesetting
\usepackage{float} % Improved interface for floating objects
\usepackage[final, colorlinks = true, 
            linkcolor = black, 
            citecolor = black]{hyperref} % For hyperlinks in the PDF
\usepackage{graphicx, multicol} % Enhanced support for graphics
\usepackage{xcolor} % Driver-independent color extensions
\usepackage{colortbl} % Para colores en tablas (Requerido por UML)
\usepackage{array} % Requerido por UML
\usepackage{marvosym, wasysym} % More symbols
\usepackage{rotating} % Rotation tools
\usepackage{censor} % Facilities for controlling restricted text
\usepackage{listings} % Environment for non-formatted code
\usepackage{pseudocode} % Environment for specifying algorithms in a natural way
\usepackage{booktabs} % Enhances quality of tables

% Configuración de listings
\lstset{
    language=Java,         
    basicstyle=\ttfamily\small,    
    keywordstyle=\color{blue!70!black}\bfseries,
    commentstyle=\color{green!50!black}\itshape,
    numbers=left,            
    frame=lines,
    literate={á}{{\'a}}1 {é}{{\'e}}1 {í}{{\'i}}1 {ó}{{\'o}}1 {ú}{{\'u}}1
             {Á}{{\'A}}1 {É}{{\'E}}1 {Í}{{\'I}}1 {Ó}{{\'O}}1 {Ú}{{\'U}}1
             {ñ}{{\~n}}1 {Ñ}{{\~N}}1 {¿}{{\?`}}1 {¡}{{\!`}}1             
}

% Configuración TikZ y Árboles
\usepackage{tikz}
\usepackage{tikz-qtree} % Easy tree drawing tool
\usetikzlibrary{arrows.meta,positioning,fit,backgrounds,calc,shadows} % Librerías para UML
\tikzset{every tree node/.style={align=center,anchor=north},
         level distance=2cm} % Configuration for q-trees

\usepackage[backend=biber,style=numeric,
            sorting=nyt]{biblatex} % Complete reimplementation of bibliographic facilities
\usepackage{csquotes} % Context sensitive quotation facilities

\usepackage[yyyymmdd]{datetime} % Uses YEAR-MONTH-DAY format for dates
\renewcommand{\dateseparator}{-} % Sets dateseparator to '-'

\usepackage{fancyhdr} % Headers and footers
\pagestyle{fancy} % All pages have headers and footers
\fancyhead{}\renewcommand{\headrulewidth}{0pt} % Blank out the default header
\fancyfoot[L]{} % Custom footer text
\fancyfoot[C]{} % Custom footer text
\fancyfoot[R]{\thepage} % Custom footer text

\newcommand{\note}[1]{\marginpar{\scriptsize \textcolor{red}{#1}}} % Enables comments in red on margin

%----------------------------------------------------------------------------------------
%	UML DIAGRAM CONFIGURATION
%----------------------------------------------------------------------------------------

% --- Color palette ---
\definecolor{umlHeader}{RGB}{215,232,252}
\definecolor{umlHeader2}{RGB}{226,239,218}
\definecolor{umlBorder}{RGB}{54,84,134}
\definecolor{namespaceFill}{RGB}{248,250,253}
\definecolor{assocColor}{RGB}{70,70,70}
\definecolor{inheritColor}{RGB}{45,78,120}

\newcolumntype{L}[1]{>{\raggedright\arraybackslash}p{#1}}

% UML class macro: #1 node name, #2 position, #3 width, #4 header color, #5 class name, #6 attributes, #7 methods
\newcommand{\umlclass}[7]{%
  \node[inner sep=0pt, outer sep=0pt, drop shadow={opacity=0.25, shadow xshift=2pt, shadow yshift=-2pt}] (#1) at #2 {%
    \scriptsize\sffamily
    \renewcommand{\arraystretch}{1.10}%
    \setlength{\tabcolsep}{4pt}%
    \begin{tabular}{|L{#3}|}
      \hline
      \rowcolor{#4}\multicolumn{1}{|c|}{\bfseries #5}\\
      \hline
      #6\\
      \hline
      #7\\
      \hline
    \end{tabular}%
  };%
}

\begin{document}

%----------------------------------------------------------------------------------------
%	TITLE SECTION
%----------------------------------------------------------------------------------------

\fancyhead[C]{}
\hrule \medskip % Upper rule
\begin{minipage}{0.295\textwidth} % Left side of title section
\raggedright
UNAL\\ % Your lecture or course
\footnotesize % Authors text size
\hfill\\
Mateo Quiceno Zapata\\
Cristian Ruiz Hernandez
\end{minipage}
\begin{minipage}{0.4\textwidth} % Center of title section
\centering 
\large % Title text size
Actividad 6\\ % Assignment title and number
\normalsize % Subtitle text size
Programación Orientada a Objetos (3007744)\\ % Assignment subtitle
\end{minipage}
\begin{minipage}{0.295\textwidth} % Right side of title section
\raggedleft
\today\\ % Date
\footnotesize % Email text size
\hfill\\
mquicenoz@unal.edu.co\\
cruizh@unal.edu.co
\end{minipage}
\medskip\hrule % Lower rule
\bigskip

\begin{center}
\Large\textbf{Entrega Actividad 6: Ejercicios Propuestos}
\end{center}
\vspace{1cm}

%----------------------------------------------------------------------------------------
%	ARTICLE CONTENTS
%----------------------------------------------------------------------------------------

\section{Ejercicio 2.10. Sobrecarga de métodos}
\textbf{Enunciado:}
Definir una clase denominada \texttt{Suma}, la cual tiene varios métodos \texttt{sumar} sobrecargados:
\begin{itemize}
    \item Un método \texttt{sumar} que obtiene la suma de dos valores enteros pasados como parámetros.
    \item Un método \texttt{sumar} que obtiene la suma de tres valores enteros pasados como parámetros.
    \item Un método \texttt{sumar} que obtiene la suma de dos valores double pasados como parámetros.
    \item Un método \texttt{sumar} que obtiene la suma de tres valores double pasados como parámetros.
\end{itemize}

\subsection{Diagrama de clases}
\begin{center}
\begin{tikzpicture}[
  font=\sffamily,
  label/.style={font=\scriptsize\sffamily, fill=white, inner sep=1.5pt}
]
\umlclass{Suma}{(0,0)}{5cm}{umlHeader}{Suma}{%
  % Atributos
}{%
  + sumar(a: int, b: int): int\\
  + sumar(a: int, b: int, c: int): int\\
  + sumar(a: double, b: double): double\\
  + sumar(a: double, b: double, c: double): double\\
  + main(args: String[])\$
}
\end{tikzpicture}
\end{center}

\subsection{Código Fuente}
\lstinputlisting[language=Java]{Ejercicio_2_10/Suma.java}


\section{Ejercicio 2.11. Sobrecarga de constructores}
\textbf{Enunciado:}
\begin{itemize}
    \item Definir una clase \texttt{Empleado} que tiene como atributos: identificador, nombre, apellidos y edad. Contiene dos constructores: uno sin parámetros con valores por defecto y otro con parámetros.
    \item Definir una clase \texttt{Caja} que tiene como atributos base, anchura, altura y tipo. Contiene varios constructores sobrecargados, incluyendo uno que invoca a otro mediante \texttt{this()}.
\end{itemize}

\subsection{Diagrama de clases}
\begin{center}
\begin{tikzpicture}[
  font=\sffamily,
  label/.style={font=\scriptsize\sffamily, fill=white, inner sep=1.5pt}
]
\umlclass{Empleado}{(0,0)}{6.5cm}{umlHeader}{Empleado}{%
  - identificador: int\\
  - nombre: String\\
  - apellidos: String\\
  - edad: int
}{%
  + Empleado()\\
  + Empleado(identificador: int, nombre: String, apellidos: String, edad: int)\\
  + imprimir(): void\\
  + main(args: String[])\$
}

\umlclass{Caja}{(8,0)}{6.5cm}{umlHeader2}{Caja}{%
  - base: double\\
  - anchura: double\\
  - altura: double\\
  - tipo: String
}{%
  + Caja(base: double, anchura: double, altura: double)\\
  + Caja()\\
  + Caja(longitud: double)\\
  + Caja(base: double, anchura: double, altura: double, tipo: String)\\
  + imprimir(): void\\
  + main(args: String[])\$
}
\end{tikzpicture}
\end{center}

\subsection{Código Fuente}
\textbf{Clase Empleado:}
\lstinputlisting[language=Java]{Ejercicio_2_11/Empleado.java}

\textbf{Clase Caja:}
\lstinputlisting[language=Java]{Ejercicio_2_11/Caja.java}


\section{Ejercicio 4.4. Polimorfismo}
\textbf{Enunciado:}
¿Cuál es el resultado de la ejecución del siguiente programa basado en el ejercicio anterior (donde se hace un casting de la instancia de la subclase a la superclase)?

\subsection{Diagrama de clases}
\begin{center}
\begin{tikzpicture}[
  font=\sffamily,
  inheritance/.style={draw=inheritColor, line width=0.8pt, -{Triangle[open, length=6pt, width=8pt]}}
]
\umlclass{Profesor}{(0,3)}{4cm}{umlHeader}{Profesor}{%
}{%
  \# imprimir(): void
}

\umlclass{ProfesorTitular}{(0,0)}{4cm}{umlHeader2}{ProfesorTitular}{%
}{%
  \# imprimir(): void
}

\draw[inheritance] (ProfesorTitular.north) -- (Profesor.south);
\end{tikzpicture}
\end{center}

\subsection{Código Fuente y Resultado}
\textbf{Clase Profesor y ProfesorTitular:}
\lstinputlisting[language=Java]{Ejercicio_4_4/Profesores/Profesor.java}
\lstinputlisting[language=Java]{Ejercicio_4_4/Profesores/ProfesorTitular.java}

\textbf{Clase Prueba:}
\lstinputlisting[language=Java]{Ejercicio_4_4/Profesores/Prueba.java}


\section{Ejercicio 4.6. Métodos polimórficos}
\textbf{Enunciado:}
¿Cuál es el resultado de la ejecución del siguiente programa que utiliza un Vector de objetos \texttt{Profesor} instanciados tanto como \texttt{Profesor} y como \texttt{ProfesorTitular}?

\subsection{Diagrama de clases}
\begin{center}
\begin{tikzpicture}[
  font=\sffamily,
  inheritance/.style={draw=inheritColor, line width=0.8pt, -{Triangle[open, length=6pt, width=8pt]}}
]
\umlclass{Profesor}{(0,3.5)}{4cm}{umlHeader}{Profesor}{%
}{%
  \# imprimir(): void
}

\umlclass{ProfesorTitular}{(0,0)}{4cm}{umlHeader2}{ProfesorTitular}{%
  \~{} años: int = 0
}{%
  \# imprimir(): void\\
  \# imprimirAños(): void
}

\draw[inheritance] (ProfesorTitular.north) -- (Profesor.south);
\end{tikzpicture}
\end{center}

\subsection{Código Fuente y Resultado}
\textbf{Clase Prueba:}
\lstinputlisting[language=Java]{Ejercicio_4_6/Profesores/Prueba.java}


\section{Ejercicio 4.7. Clases abstractas}
\textbf{Enunciado:}
\begin{itemize}
    \item Definir una clase abstracta denominada \texttt{Numérica} con métodos abstractos: \texttt{toString()}, \texttt{equals()}, \texttt{sumar()}, \texttt{restar()}, \texttt{multiplicar()} y \texttt{dividir()}.
    \item Definir una clase \texttt{Fracción} que representa un número fraccionario, hereda de \texttt{Numérica}, implementa sus métodos y tiene numerador y denominador.
    \item Crear una clase de prueba que utilice los métodos implementados.
\end{itemize}

\subsection{Diagrama de clases}
\begin{center}
\begin{tikzpicture}[
  font=\sffamily,
  inheritance/.style={draw=inheritColor, line width=0.8pt, -{Triangle[open, length=6pt, width=8pt]}}
]
\umlclass{Numerica}{(0,4.5)}{5cm}{umlHeader}{\textit{Numérica}}{%
}{%
  + toString(): String \textit{abstract}\\
  + equals(ob: Object): boolean \textit{abstract}\\
  + sumar(numero: Numérica): Numérica \textit{abstract}\\
  + restar(numero: Numérica): Numérica \textit{abstract}\\
  + multiplicar(numero: Numérica): Numérica \textit{abstract}\\
  + dividir(numero: Numérica): Numérica \textit{abstract}
}

\umlclass{Fraccion}{(0,0)}{5.5cm}{umlHeader2}{Fracción}{%
  - numerador: int\\
  - denominador: int
}{%
  + Fracción(numerador: int, denominador: int)\\
  + getNumerador(): int\\
  + getDenominador(): int\\
  + toString(): String\\
  + equals(ob: Object): boolean\\
  + sumar(numero: Numérica): Numérica\\
  + restar(numero: Numérica): Numérica\\
  + multiplicar(numero: Numérica): Numérica\\
  + dividir(numero: Numérica): Numérica\\
  - simplificar(f: Fracción): Fracción\\
  - mcd(a: int, b: int): int
}

\draw[inheritance] (Fraccion.north) -- (Numerica.south);
\end{tikzpicture}
\end{center}

\subsection{Código Fuente}
\textbf{Clase Numérica:}
\lstinputlisting[language=Java]{Ejercicio_4_7/Numerica.java}

\textbf{Clase Fracción:}
\lstinputlisting[language=Java]{Ejercicio_4_7/Fraccion.java}

\textbf{Clase Prueba:}
\lstinputlisting[language=Java]{Ejercicio_4_7/Prueba.java}

\end{document}
